\section{Background}\label{sec:background}

\subsection{$\Sigma$-protocols}\label{sigma_protocol}
A $\Sigma$-protocol~\cite{cramer1996modular,cramer1994proofs} is a
\rocq[three-move proof]{sigma_proto} for an NP relation $R$. A prover
wants to convince a verifier that it knows a secret \emph{witness} $w$
for a public statement $x$. The prover sends an announcement $a$, the
verifier replies with a random challenge $c$, the prover sends a response
$r$, and the verifier accepts or rejects the transcript $(a,c,r)$. The
protocol is \rocq[\emph{complete}]{schnorr_completeness} if an honest prover is always accepted.
It is \rocq[\emph{special sound}]{special_soundness_berry} if a witness can be computed from any two
accepting transcripts $(a,c,r)$ and $(a,c',r')$ with $c \ne c'$, so a
prover who can answer two challenges must know a witness. It is
\rocq[\emph{special honest-verifier zero-knowledge}]{special_honest_verifier_zkp_perm} (SHVZK) if, for every
challenge $c$, a simulator that does not know the witness produces
transcripts with the same distribution as the honest prover, so a
transcript reveals nothing about the witness. The Fiat--Shamir
transformation removes the interaction: the prover computes $c$ itself,
as the hash of the statement and the announcement. The best-known example is
Schnorr's \rocq[proof of knowledge of a discrete logarithm]{schnorr_protocol}~\cite{schnorr1991efficient}.
Deployed voting systems such as Helios~\cite{adida2008helios},
Belenios~\cite{belenios2019}, and the Swiss Post
system~\cite{10.1007/978-3-031-15911-4_4} build their proofs from Schnorr
proofs and their compositions: conjunctions (``and''), disjunctions
(``or''), and proofs that two discrete logarithms are equal
(Chaum--Pedersen). Applied to ElGamal ciphertexts, these proofs show that
a ciphertext encrypts one of a list of messages, or that it decrypts to
a given value. Every proof of the scheme in Section~\ref{sec:schulze} is
of this kind.

\subsection{The Schulze method}\label{sec:bg:schulze}
A ranked ballot orders the $m$ candidates, with ties allowed. The
natural winner of a ranked election is a \emph{Condorcet winner}: a
candidate who beats every other candidate head to head. There may be
none, because the head-to-head contests can form a cycle or end in ties.
The Schulze method~\cite{schulze2011} always produces a winner, and it
elects the Condorcet winner whenever there is one. It uses the ballots
only through their pairwise counts.

The \rocq[\emph{pairwise count matrix}]{margin_count} $D$ (\emph{count matrix} for short) records
pairwise preferences: $D_{ij}$ is the number of voters who rank $i$
strictly above $j$, so a voter who ties $i$ and $j$ adds to neither
$D_{ij}$ nor $D_{ji}$. Candidate $i$ \rocq[\emph{defeats}]{defeat} $j$ with strength
$d_{ij} = D_{ij}$ if $D_{ij} > D_{ji}$, and $d_{ij} = 0$ otherwise. We
call $d$ the \emph{defeat matrix}. This is the \emph{winning-votes}
measure of defeat strength. 

Candidate $i$ may beat $j$ indirectly, through other candidates, without
defeating it. A path $i = u_0, \ldots, u_\ell = j$ is a chain of defeats,
and it is only \rocq[as strong as its weakest one]{path_strength},
$\min_t d_{u_t u_{t+1}}$. The strongest-path value $p_{ij}$ is the
strength of the best path from $i$ to $j$. Candidate $w$ is a
\rocq[\emph{Schulze winner}]{schulze_winner} if $p_{wj} \ge p_{jw}$ for every $j$: no candidate
has a stronger path to $w$ than $w$ has back. At least one Schulze
winner always exists: Schulze shows that the relation $p_{ij} > p_{ji}$
is transitive~\cite{schulze2011}, so some candidate is beaten by no one
in this sense. A Condorcet winner $w$ is the unique one: it
defeats every $j$, so $p_{wj} \ge d_{wj} > 0$, and nobody defeats $w$, so
every path into $w$ ends in an edge of strength $0$ and $p_{jw} = 0$.

\textbf{Running example:} Five voters rank three candidates:
\[
  \begin{array}{c}
    A > B = C, \qquad B > C > A, \qquad B > C > A,\\
    C > A > B, \qquad A > C > B.
  \end{array}
\]
Their pairwise counts, defeats, and strongest paths are
\[
  \setlength{\arraycolsep}{2.5pt}
  D =\begin{array}{c|ccc}
        & A & B & C \\ \hline
      A & - & 3 & 2 \\
      B & 2 & - & 2 \\
      C & 3 & 2 & -
      \end{array}
  \quad
  d = \begin{array}{c|ccc}
        & A & B & C \\ \hline
      A & - & 3 & 0 \\
      B & 0 & - & 0 \\
      C & 3 & 0 & -
      \end{array}
  \quad
  p = \begin{array}{c|ccc}
        & A & B & C \\ \hline
      A & - & 3 & 0 \\
      B & 0 & - & 0 \\
      C & 3 & 3 & -
      \end{array}
\]
(rows are $i$, columns $j$). For instance, $D_{AB} = 3$ because ballots
$1$, $4$, and $5$ rank $A$ above $B$. $D_{BC} = D_{CB} = 2$ because
ballots $2$ and $3$ rank $B$ above $C$, ballots $4$ and $5$ rank $C$
above $B$, and ballot $1$ ties them and counts for neither. There are two defeats, $A$ over $B$ and $C$
over $A$, each of strength $3$; $B$ and $C$ tie $2$ to $2$. No candidate
beats both others, so there is no Condorcet winner. $C$ has no direct
defeat of $B$, but the path $C \to A \to B$ has weakest link $3$, so
$p_{CB} = 3$. Every other entry of $p$ is a single defeat or $0$. $C$ is
the unique Schulze winner, since $p_{CA} = 3 \ge p_{AC} = 0$ and $p_{CB}
= 3 \ge p_{BC} = 0$, while $A$ fails against $C$ ($p_{AC} = 0 < 3$) and
$B$ fails against both.

In algebraic terms, the matrix $p$ is the \rocq[closure]{is_closure} of
$d$ in the semiring with $\max$ as addition and $\min$ as multiplication:
a path's strength is the $\min$ of its edges, and $p_{ij}$ is the $\max$
over paths. We call it the \emph{min-max closure}. It is
\rocq[computable]{minmax_closure} in $O(m^3)$ operations by
Floyd--Warshall or Kleene
iteration~\cite{10.1007/978-3-319-66107-0_26}. Floyd--Warshall
starts from $p = d$ and, for each candidate $k$ in turn, sets $p_{ij}
\leftarrow \max(p_{ij}, \min(p_{ik}, p_{kj}))$: a path through $k$
replaces the current best if its weaker half is stronger. In the running
example the step $k = A$ sets $p_{CB} = \max(0, \min(d_{CA}, d_{AB})) =
3$, and no other step changes an entry. \rocq[Three properties]{cert_ok}
identify the closure, and our certificate proves them over encrypted
values (Section~\ref{sec:tally}): $p_{ij} \ge d_{ij}$; $p_{ij} \ge
\min(p_{ik}, p_{kj})$ for every $k$; and $p_{ij}$ is the strength of some
actual path. The winner depends only on $D$, and $D$ is a sum of
per-ballot bits. So the winner can be computed from the homomorphic sum
of the encrypted ballots, and no individual ballot is ever decrypted.

The Schulze method is used by Debian, the Free Software Foundation
Europe, and other organisations that vote online with a few hundred to a
few thousand voters.

\subsection{Tally hiding}\label{sec:bg:hiding}
A verifiable ranked tally that publishes the ballots allows the
\emph{Italian attack} of
Section~\ref{sec:intro}~\cite{teague2008stv,shufflesum2009}. Publishing
only $D$ leaks less, but it still leaks in a small electorate.
Following Kryvos~\cite{kryvos2022}, a system is \emph{fully}
tally-hiding if nobody, not even the talliers, learns anything beyond the
result. It is \emph{publicly} tally-hiding if the talliers learn the
aggregated tally (here $D$) and the public learns only the result. The scheme of this paper is publicly tally-hiding except for one leak:
the witness paths of its certificate, which reveal who wins some
head-to-head contests but not by how much (Section~\ref{sec:security}).
